\documentclass[12pt]{article}

\usepackage[margin=1.12in]{geometry}
\usepackage{amsmath,amssymb,amsthm,mathtools}
\usepackage{array}
\usepackage{booktabs}
\usepackage{hyperref}

\newtheorem{theorem}{Theorem}
\newtheorem{proposition}{Proposition}
\theoremstyle{remark}
\newtheorem{remark}{Remark}

\newcommand{\A}{\mathcal A}
\newcommand{\C}{\mathbb C}
\newcommand{\M}{\mathcal M}
\newcommand{\Ocal}{\mathcal O}
\newcommand{\rce}{\operatorname{rce}}

\setlength{\emergencystretch}{3em}
\hypersetup{
  pdftitle={A Constant Finite-Center LCQFT Extension and the Accountability Premise},
  pdfsubject={Locally covariant QFT, dynamical locality, finite centers, and quotient-visible data},
  pdfkeywords={LCQFT, dynamical locality, relative Cauchy evolution, finite centers, accountability}
}

\title{A Constant Finite-Center LCQFT Extension\\
and the Accountability Premise}
\author{}
\date{Live review note of June 20, 2026}

\begin{document}
\maketitle

\begin{abstract}
Let \(\A\) be a BFV locally covariant quantum field theory valued in unital
\(*\)-algebras, and let \(Z_q=\C^q\).  The constant finite-center extension
\[
  \A_q(M)=\A(M)\otimes Z_q,\qquad
  \A_q(\psi)=\A(\psi)\otimes \operatorname{id}_{Z_q}
\]
preserves functoriality, injective morphisms, and the time-slice axiom.  Its
relative Cauchy evolution is
\[
  \rce_q[h]=\rce[h]\otimes \operatorname{id}_{Z_q}.
\]
If dynamical locality is understood in the Fewster-Verch RCE sense, and the
target tensor product identifies \(\A(M)\otimes \C^q\) with
\(\bigoplus_{j=1}^q\A(M)\), then dynamical locality is inherited from
\(\A\): both the kinematic and dynamical local subobjects are tensored by
\(Z_q\).  Nevertheless any record that factors through a character
\(\pi_j:Z_q\to \C\) is blind to the nonselected finite-center directions.  Thus
BFV functoriality, time-slice, RCE, and RCE dynamical locality alone do not
reconstruct hidden finite central rank.  A positive result must add the exact
premise that removes or measures the center: extended locality, local
definiteness, center reduction, indecomposability, or a center-separating
measurement channel.
\end{abstract}

\paragraph{Status.}
This is non-frozen external-review metadata for Team 1.  It is intended as a
small claim that an LCQFT referee can classify as known, false,
missing-hypothesis, or useful.  It is not a new construction of LCQFT, and it
does not modify any frozen Team 1 manuscript, Lean archive, public mirror
manifest, or SHA-256 package.

\section{Setup}

Let \(\mathsf{Loc}\) denote the usual BFV category of globally hyperbolic
spacetimes and admissible embeddings, and let \(\mathsf{Alg}\) be a category of
unital \(*\)-algebras with injective morphisms.  A locally covariant theory is
a functor
\[
  \A:\mathsf{Loc}\longrightarrow \mathsf{Alg}.
\]
It obeys the time-slice axiom if \(\A(\psi)\) is an isomorphism whenever
\(\psi\) is a Cauchy morphism.  Relative Cauchy evolution is the automorphism
of \(\A(M)\) obtained from the time-slice zig-zag associated with a compactly
supported metric perturbation.

Fix \(q\geq 2\), let \(Z_q=\C^q\), and define
\[
  \A_q(M)=\A(M)\otimes Z_q,\qquad
  \A_q(\psi)=\A(\psi)\otimes \operatorname{id}_{Z_q}.
\]
For a character \(\pi_j:Z_q\to\C\), define the quotient-visible natural map
\[
  Q_{j,M}=\operatorname{id}_{\A(M)}\otimes \pi_j:\A_q(M)\to \A(M).
\]

\begin{theorem}[constant finite-center extension]
\label{thm:constant-center}
If \(\A\) is a BFV locally covariant theory with injective morphisms, then
\(\A_q\) is also a BFV locally covariant theory with injective morphisms.  If
\(\A\) obeys the time-slice axiom, then so does \(\A_q\), and
\[
  \rce_q[h]=\rce[h]\otimes \operatorname{id}_{Z_q}.
\]
Any operational, sector, recovery, modular, or RCE record factoring through
\(Q_{j,M}\) is insensitive to the nonselected central idempotents and to every
nonconstant central contrast in \(\ker\pi_j\).
\end{theorem}

\begin{proof}
Functoriality, identities, and composition are componentwise.  Injectivity is
preserved because tensoring with the finite-dimensional algebra \(Z_q\) is
exact here; equivalently, \(\A(M)\otimes \C^q\simeq\bigoplus_{j=1}^q\A(M)\),
and \(\A(\psi)\otimes\operatorname{id}_{Z_q}\) is the direct sum of \(q\)
copies of \(\A(\psi)\).  If \(\psi\) is Cauchy, \(\A(\psi)\) is an isomorphism,
so its finite direct sum is an isomorphism.  Thus the time-slice axiom is
preserved.  The relative Cauchy evolution in \(\A_q\) is obtained by applying
the same time-slice zig-zag with every arrow tensored by
\(\operatorname{id}_{Z_q}\), giving the stated formula.

Finally, if \(R_q=R\circ Q_{j,M}\), then for \(a\in\A(M)\) and
\(z,z'\in Z_q\) with \(\pi_j(z)=\pi_j(z')\),
\[
  R_q(a\otimes z)=R(a\,\pi_j(z))=R(a\,\pi_j(z'))=R_q(a\otimes z').
\]
Thus the record is constant on the hidden finite-center fiber.
\end{proof}

\section{Dynamical Locality}

Fewster-Verch dynamical locality compares the kinematic local subobject
\[
  \A_{\mathrm{kin}}(M;O)=\A(\iota_O)(\A(M|_O))
\]
with the dynamical subobject
\[
  \A_{\mathrm{dyn}}(M;O)
  =
  \{a\in \A(M): \rce_M[h](a)=a
  \text{ for all } h \text{ supported in } O^\perp\}.
\]
The theory is dynamically local when these two subobjects agree for all
admissible \(O\).

\begin{proposition}[inheritance of RCE dynamical locality]
\label{prop:dynloc}
Assume \(\A\) is dynamically local in the Fewster-Verch RCE sense, and assume
the tensor product in the target algebra category identifies
\(\A(M)\otimes\C^q\) with \(\bigoplus_{j=1}^q\A(M)\).  Then the constant
finite-center extension \(\A_q\) is dynamically local in the same sense.
\end{proposition}

\begin{proof}
For every admissible \(O\subset M\),
\[
  (\A_q)_{\mathrm{kin}}(M;O)
  =
  \A_{\mathrm{kin}}(M;O)\otimes \C^q
\]
by functoriality of the inclusion \(\iota_O:M|_O\to M\).  By
Theorem~\ref{thm:constant-center},
\[
  \rce_q[h]=\rce[h]\otimes \operatorname{id}_{\C^q}.
\]
Under the finite direct-sum identification, an element of \(\A_q(M)\) is a
tuple \((a_1,\ldots,a_q)\).  It is fixed by all \(\rce_q[h]\) with
\(\operatorname{supp}h\subset O^\perp\) exactly when every \(a_j\) is fixed by
the corresponding \(\rce[h]\).  Hence
\[
  (\A_q)_{\mathrm{dyn}}(M;O)
  =
  \A_{\mathrm{dyn}}(M;O)\otimes \C^q.
\]
If \(\A_{\mathrm{kin}}(M;O)=\A_{\mathrm{dyn}}(M;O)\), then
\((\A_q)_{\mathrm{kin}}(M;O)=(\A_q)_{\mathrm{dyn}}(M;O)\).
\end{proof}

\begin{proposition}[extended locality and local definiteness are positive exits]
\label{prop:extended-locality}
Let \(q>1\).  Suppose a spacetime \(M\) contains two nonempty causally disjoint
admissible regions \(O_1,O_2\), and let their kinematic images in
\(\A_q(M)\) be denoted by \((\A_q)_{\mathrm{kin}}(M;O_i)\).  Then
\[
  \C 1_{\A(M)}\otimes Z_q
  \subset
  (\A_q)_{\mathrm{kin}}(M;O_1)
  \cap
  (\A_q)_{\mathrm{kin}}(M;O_2).
\]
Consequently \(\A_q\) violates any extended-locality axiom requiring causally
disjoint nonempty local algebras to intersect only in scalars.  Likewise, if a
local-definiteness axiom requires the intersection of algebras for regions
shrinking to a point to be only \(\C 1\), then \(\A_q\) violates that axiom for
the same reason: the constant finite center remains present in every local
kinematic algebra.
\end{proposition}

\begin{proof}
For every admissible inclusion \(\iota_O:M|_O\to M\), the element
\(1_{\A(M|_O)}\otimes z\in \A(M|_O)\otimes Z_q\) is mapped by
\(\A(\iota_O)\otimes\operatorname{id}_{Z_q}\) to
\(1_{\A(M)}\otimes z\).  Hence the same central copy
\(\C 1_{\A(M)}\otimes Z_q\) lies in the kinematic image of every nonempty
region.  If \(q>1\), this copy contains non-scalar central idempotents, so the
intersection of two causally disjoint local algebras is not just \(\C 1\).
The shrinking-region statement is identical: every algebra in the shrinking
family contains the same central copy.
\end{proof}

\section{Positive Exits}

Proposition~\ref{prop:dynloc} is narrow.  It says that the RCE-defined equality
of kinematic and dynamical local subobjects does not by itself remove a
constant finite center.  Proposition~\ref{prop:extended-locality} gives the
matching positive exit: stronger physical assumptions can remove it.

\begin{center}
\begin{tabular}{>{\raggedright\arraybackslash}p{0.30\linewidth}>{\raggedright\arraybackslash}p{0.56\linewidth}}
\toprule
Exit premise & Effect on \(\A_q\) for \(q>1\) \\
\midrule
Extended locality & If algebras of causally disjoint nonempty regions may
intersect only in scalars, \(\A_q\) fails because every local algebra contains
the common central copy of \(Z_q\). \\
\addlinespace
Local definiteness & If the intersection of local algebras shrinking to a
point must be \(\C\), \(\A_q\) fails because \(Z_q\) remains in every local
algebra. \\
\addlinespace
Center reduction & If the framework quotients or reduces away central
degrees of freedom, as in reduced Maxwell-type constructions, the missing
physical-exclusion premise has been supplied. \\
\addlinespace
Center-separating measurement & If the probe/readout is tomographically
complete on \(Z_q\), the data do not factor through a single character
\(\pi_j\); the quotient-visible hypothesis is false. \\
\bottomrule
\end{tabular}
\end{center}

\section{Referee Question}

The classification question is precise:

\begin{quote}
Is the constant finite-center extension above already excluded by the standard
LCQFT hypotheses intended in the target setting?  If so, name the axiom or
theorem.  If not, BFV functoriality, time-slice, RCE, and RCE dynamical
locality alone do not supply finite-center accountability.
\end{quote}

The desired external response is not encouragement.  It is one of:
known theorem, false statement, missing positive hypothesis, physically too
broad, or useful expository no-go.

\begin{thebibliography}{99}

\bibitem{BFV}
R. Brunetti, K. Fredenhagen, and R. Verch,
The generally covariant locality principle: a new paradigm for local quantum
field theory,
\emph{Communications in Mathematical Physics} \textbf{237} (2003), 31--68;
\texttt{arXiv:math-ph/0112041}.

\bibitem{FewsterVerch}
C. J. Fewster and R. Verch,
Dynamical locality and covariance: what makes a physical theory the same in
all spacetimes?,
\emph{Annales Henri Poincare} \textbf{13} (2012), 1613--1674;
\texttt{arXiv:1106.4785}.

\bibitem{FewsterOverview}
C. J. Fewster,
The paradigm of local covariance,
in \emph{Quantum Field Theory and Gravity}, Birkhauser, 2012;
see also \texttt{arXiv:1105.6202}.

\bibitem{FewsterVerchMeasurement}
C. J. Fewster and R. Verch,
Quantum fields and local measurements,
\emph{Communications in Mathematical Physics} \textbf{378} (2020), 851--889;
\texttt{arXiv:1810.06512}.

\bibitem{FewsterVerchReview}
C. J. Fewster and R. Verch,
Measurement in quantum field theory,
\emph{Lecture Notes in Physics} \textbf{1017} (2024), 1--45;
\texttt{arXiv:2304.13356}.

\bibitem{FewsterLang}
C. J. Fewster and B. Lang,
Dynamical locality of the free Maxwell field,
\emph{Annales Henri Poincare} \textbf{17} (2016), 401--436;
\texttt{arXiv:1403.7083}.

\end{thebibliography}

\end{document}
